% ------------------------------------------------------------------------
% file `2_noel_2021_part2-exercise-4-exercise-body.tex'
%   in folder `xsim/'
%
%     exercise of type `exercise' with id `4'
%
% generated by the `exercise' environment of the
%   `xsim' package v0.20c (2021/02/03)
% from source `2_noel_2021_part2' on 2021/12/10 on line 62
% ------------------------------------------------------------------------
    Coche la réponse correcte.

    \begin{enumerate}
        \item $2n+1$ et $2n+2$ sont deux nombres \dots
              \begin{qcm}[cols=4]
                  \item pairs
                  \item impairs
                  \item consécutifs
                  \item premiers
              \end{qcm}
        \item Dans l'égalité $49=4\cdot10+9$, le nombre $10$ est \dots
              \begin{qcm}[cols=4]
                  \item le diviseur
                  \item le dividende
                  \item le quotient
                  \item le reste
              \end{qcm}
        \item $(-2)^3$ est égale à \dots
              \begin{qcm}[cols=4]
                  \item $8$
                  \item $6$
                  \item $-8$
                  \item $-6$
              \end{qcm}
        \item $\left(3x^2\right)^2$ se réduit en \dots
              \begin{qcm}[cols=4]
                  \item $3x^4$
                  \item $9x^4$
                  \item $9x^2$
                  \item $6x^4$
              \end{qcm}
    \end{enumerate}
